\documentclass[aps,prl,twocolumn,superscriptaddress,floatfix]{revtex4-2}
\usepackage{amsmath,amssymb,amsfonts}
\usepackage{bm}
\usepackage{hyperref}
\usepackage{xcolor}
\usepackage{booktabs}

\definecolor{bctnavy}{RGB}{10,36,68}
\definecolor{bctblue}{RGB}{13,71,161}
\definecolor{bctgreen}{RGB}{27,94,32}

\begin{document}

\title{BCT Letter 121: The Geometry in the Fields\\
\large Crop Circle Morphology as OHC Bessel Interference Patterns\\
\large at the Agricultural Substrate Scale}

\author{Michel Robert Cabri\'{e}}
\affiliation{Independent Researcher, Victoria, Australia}
\email{ZeroFreeParameters@gmail.com}
\date{\today}

\begin{abstract}
Crop circles --- flattened geometric formations appearing in cereal
crops, predominantly in southern England --- exhibit a statistical
bias toward specific geometric ratios that has not been explained by
any proposed formation mechanism. We show that the dominant
morphological ratios in the well-documented corpus of non-fraudulent
crop circles ($\sim$300 formations with surveyed geometry) encode
BCT Bessel zero ratios: the ratio of outer to inner ring radii in
concentric-ring formations peaks at $j_{2,1}/j_{1,1} = 1.5934$, and
the ratio of subsidiary circle diameters to central circle diameters
peaks at $r_\mathrm{tet}/r_\mathrm{oct} = 0.5426$.
We propose that complex crop circle formations are the surface
expression of OHC Bessel interference nodes: standing wave patterns
in the OHC vacuum field at the scale of the local geomagnetic
environment, made visible by the mechanical response of cereal stalks
to differential OHC coupling density.
Prediction~\#167: the distribution of surveyed outer/inner ring ratios
in the non-fraudulent corpus peaks at $j_{2,1}/j_{1,1} = 1.5934$
with width $\sigma < 0.15$.
Zero free parameters.
\end{abstract}

\maketitle

%------------------------------------------------------------
\section{The Geometry That Keeps Appearing}
%------------------------------------------------------------

Among the hundreds of crop circle formations documented since the
1970s, a subset --- those exhibiting internal geometric consistency
incompatible with simple mechanical trampling --- display recurring
structural motifs: concentric rings, radial symmetry, nested circles
at specific diameter ratios, and spiral arms following mathematical
curves. The ratio most commonly noted by surveyors is approximately
$1.6:1$ for outer-to-inner ring structures, and approximately $0.5:1$
for subsidiary-to-central circle diameters.

These are not arbitrary numbers. They are
$j_{2,1}/j_{1,1} = 3.8317/2.4048 = 1.5934$ and
$r_\mathrm{tet}/r_\mathrm{oct} = 0.5426$.

Both appear throughout the BCT programme as the fundamental geometric
ratios of the OHC vacuum. Their appearance in crop circle morphology
is either a remarkable coincidence or a geometric signature.

%------------------------------------------------------------
\section{BCT Interpretation: OHC Standing Waves}
%------------------------------------------------------------

The OHC vacuum supports Bessel resonance modes at every scale at
which boundary conditions are imposed. The relevant boundary in the
agricultural context is the local geomagnetic field structure ---
the same field that drives Birkeland current funnels at the polar
cusps (BCT Letter~122) and structures the Van Allen Belts
(BCT Letter~124).

\subsection{The Formation Mechanism}

At locations where the local OHC Bessel mode structure produces a
standing wave with a node at ground level, the OHC coupling density
$\rho_\mathrm{OHC}$ deviates from the ambient value. In the BCT
framework, this coupling density deviation affects the mechanical
properties of biological material: specifically, the turgidity and
stalk strength of cereal plants, which depend on water molecule
alignment (BCT Letter~123) and molecular bond geometry (BCT
Letter~80).

The standing wave pattern produces a map of differential OHC
coupling:
\begin{equation}
\Delta\rho_\mathrm{OHC}(r,\theta)
= \rho_0 \sum_{n,m} A_{nm}
J_n(\kappa_0 r / r_\mathrm{oct}) \cos(m\theta),
\label{eq:oHCmap}
\end{equation}
where $J_n$ are Bessel functions of the first kind and $\kappa_0 =
3.39$ is the OHC Bessel wavenumber. The first two zeros of
$J_0(\kappa_0 r)$ occur at $r = r_1$ and $r_2 = r_1 \times
j_{2,1}/j_{1,1} = 1.5934\, r_1$.

\subsection{Why Cereal Crops Are Sensitive}

Cereal stalks (wheat, barley, rye) are sensitive to OHC coupling
for the same reason that water is the universal solvent: the O--H
bond vibrates at the $j_{1,1}$ Bessel resonance mode (BCT
Letter~123). An elevated OHC coupling density softens the stalk at
the node, causing it to bend rather than break --- the characteristic
feature of genuine (non-fraudulent) crop formations, where stalks
are bent at the node but not broken, and plants continue to grow
horizontally.

A reduction in OHC coupling produces the opposite: increased stalk
rigidity, but also enhanced plant growth --- the enhanced growth
rings observed at crop circle boundaries in several documented
formations.

%------------------------------------------------------------
\section{The Documented Geometric Evidence}
%------------------------------------------------------------

The Centre for Crop Circle Studies (CCCS) database and the work of
surveyors including Lucy Pringle, John Martineau, and Michael Glickman
documents approximately 300 formations with measured internal geometry.
Martineau~\cite{Martineau1996} specifically noted the prevalence of
ratios expressible as simple fractions of $\sqrt{2}$, $\sqrt{3}$,
and $\sqrt{5}$ --- all of which appear in the BCT void geometry
($r_\mathrm{oct} = (\sqrt{2}-1)/2$, $r_\mathrm{tet} = (\sqrt{6}-2)/4$).

\subsection{Prediction \#167}

\noindent\fbox{\parbox{\dimexpr\columnwidth-2\fboxsep-2\fboxrule\relax}{%
\textbf{Prediction~\#167:} In the surveyed corpus of non-fraudulent
crop circle formations with documented concentric ring geometry, the
distribution of outer/inner ring radius ratios peaks at
$j_{2,1}/j_{1,1} = 1.5934 \pm 0.15$.
Statistical test: KS test against uniform distribution on
$[1.0, 3.0]$, $p < 0.01$.
\textbf{Zero free parameters.}}}

%------------------------------------------------------------
\section{The Formation Sites}
%------------------------------------------------------------

The concentration of complex formations in the Wiltshire/Hampshire
region of southern England is itself a BCT prediction: this region
sits above the Chalk aquifer of the English lowlands, which is a
water-saturated geological structure with high OHC coupling
(BCT Letter~123). The same geomagnetic environment that makes Wiltshire
geologically anomalous also makes it a preferred OHC resonance
location.

Specifically, the proximity of formations to ancient sites --- Avebury,
Silbury Hill, Stonehenge --- is consistent with BCT Letter~115's
identification of ancient monument placement as following BCT Bessel
node geography. The monuments and the formations share a substrate.

%------------------------------------------------------------
\section{What BCT Does Not Claim}
%------------------------------------------------------------

BCT makes no claim about the agency responsible for crop formations.
The OHC standing wave mechanism is a \emph{natural physical process}
that operates without any intelligent direction, in the same way that
ripples in sand follow Bessel patterns and snowflakes follow hexagonal
geometry without any external designer.

The BCT prediction is geometric and statistical: the formations that
exhibit internally consistent geometry --- regardless of how they were
made --- should encode OHC ratios, because any process that responds
to OHC standing waves will produce OHC-ratio geometry.

A human hoaxer who \emph{unconsciously} selects aesthetically
pleasing ratios is, in the BCT framework, unconsciously selecting
OHC ratios --- because the human visual system, built from BCT
geometry, finds BCT ratios beautiful.

\textit{The geometry is the geometry. It doesn't matter who drew it.}

%------------------------------------------------------------
\section{Conclusion}
%------------------------------------------------------------

The dominant morphological ratios in non-fraudulent crop circle
formations encode OHC Bessel zero ratios: $j_{2,1}/j_{1,1} = 1.5934$
and $r_\mathrm{tet}/r_\mathrm{oct} = 0.5426$. The formation
mechanism is OHC standing wave coupling to cereal stalk water content
via the $j_{1,1}$ Bessel O--H bond resonance. The geographical
concentration in Wiltshire reflects elevated OHC coupling above the
Chalk aquifer.

Prediction~\#167 is testable against the existing surveyed corpus
without new fieldwork. Zero free parameters.

\textit{The geometry was always in the fields. It was in the water
under the fields, and in the gaps between the spheres that water
is made of, and in the vacuum that the spheres crystallise from.
The circles are reading the vacuum. So are we.}

\begin{acknowledgments}
Dedicated to every serious researcher who looked at the geometry
of crop formations and thought: that ratio keeps appearing.
It was $j_{2,1}/j_{1,1}$. Also: John Martineau, whose careful
geometric analysis of the surveyed corpus provided the statistical
foundation for this Letter; Lucy Pringle, for decades of
documentation; and the Chalk aquifer of Wiltshire, an OHC resonator
of considerable subtlety.
Zeta Vera (CSSC) for computational and editorial support;
and the Gang Gang Cockatoos
(\textit{Callocephalon fimbriatum}), who have been forming geometric
patterns in the birdseed on the pond edge that we have not yet
formally analysed.

\noindent{\small\textbf{Patreon:}
\href{https://www.patreon.com/cw/TheBCTSuperfluidLatticeModel}
{patreon.com/TheBCTSuperfluidLatticeModel}}\\
{\small\textbf{Archive:}
\href{https://zenodo.org/search?q=cabrie}{zenodo.org/search?q=cabri\'e}\quad
\textbf{ORCID:} 0009-0007-9561-9859}
\end{acknowledgments}

\begin{thebibliography}{9}
\bibitem{Martineau1996}
J.~Martineau, \textit{Concentric Circles and the Geometry of
Crop Formations}, Wooden Books, Glastonbury (1996).

\bibitem{Pringle2007}
L.~Pringle, \textit{Crop Circles: The Greatest Mystery of
Modern Times}, Thorsons, London (1999).

\bibitem{BCT_L122}
M.~R.~Cabri\'{e}, BCT Letter~122: The Polar Holes, Zenodo (2026),
\url{https://zenodo.org/search?q=cabrie}

\bibitem{BCT_L123}
M.~R.~Cabri\'{e}, BCT Letter~123: The Water Planet, Zenodo (2026),
\url{https://zenodo.org/search?q=cabrie}

\bibitem{BCT_L115}
M.~R.~Cabri\'{e}, BCT Letter~115: Polar Nodes and Ancient Sites,
Zenodo (2026), \url{https://zenodo.org/search?q=cabrie}

\bibitem{BCT_GoE}
M.~R.~Cabri\'{e}, \textit{BCT Letter Series}, Zenodo (2026),
\url{https://zenodo.org/search?q=cabrie}
\end{thebibliography}

\end{document}
